\clearpage
\begin{explanationbox}
\textbf{\textcolor{CExplanation}{一、为什么要有这条结论？}}
\begin{center}\begin{tikzpicture}[font=\small,>=stealth,line join=round]
\draw[gray!40,rounded corners=3pt] (0,0) rectangle (14.2,8.5);\node at (7.1,7.9) {先看内层：合并同一个数，还是得到一对相反数？};
\node at (2.4,6.9) {$x^2+3$：内偶};\begin{scope}[shift={(2.4,3.8000000000000003)},x=1.1333333333333333cm,y=0.4cm]\draw[->,gray] (-1.5,0)--(1.5,0) node[right]{$x$};\draw[->,gray] (0,-3)--(0,6) node[above]{$y$};\draw[blue,thick,solid,domain=-1.3:1.3,samples=160] plot(\x,{(\x)^2+3});\fill[orange!85!black] (-1,4) circle(2pt);\fill[orange!85!black] (1,4) circle(2pt);\end{scope}\node[align=center,text width=4cm] at (2.4,1.7) {$F(-1)=F(1)=4$};\node at (7.050000000000001,6.9) {$x^3$：内奇外奇};\begin{scope}[shift={(7.050000000000001,3.8000000000000003)},x=1.1333333333333333cm,y=0.4cm]\draw[->,gray] (-1.5,0)--(1.5,0) node[right]{$x$};\draw[->,gray] (0,-3)--(0,6) node[above]{$y$};\draw[violet,thick,solid,domain=-1.3:1.3,samples=160] plot(\x,{(\x)^3});\fill[orange!85!black] (-1,-1) circle(2pt);\fill[orange!85!black] (1,1) circle(2pt);\end{scope}\node[align=center,text width=4cm] at (7.050000000000001,1.7) {$F(-1)=-1,\ F(1)=1$};\node at (11.7,6.9) {$x^2$：内奇外偶};\begin{scope}[shift={(11.7,3.8000000000000003)},x=1.1333333333333333cm,y=0.4cm]\draw[->,gray] (-1.5,0)--(1.5,0) node[right]{$x$};\draw[->,gray] (0,-3)--(0,6) node[above]{$y$};\draw[orange!85!black,thick,solid,domain=-1.3:1.3,samples=160] plot(\x,{(\x)^2});\fill[orange!85!black] (-1,1) circle(2pt);\fill[orange!85!black] (1,1) circle(2pt);\end{scope}\node[align=center,text width=4cm] at (11.7,1.7) {$F(-1)=F(1)=1$};
\node[text width=12.6cm,align=center] at (7.1,.65) {三图同尺度；分别取 $g=x^2,f(u)=u+3$，及 $g=x,f(u)=u^3$ 或 $u^2$。};
\end{tikzpicture}\end{center}


我们判断偶函数，本来要比较 $F(2)$ 和 $F(-2)$；判断奇函数，则要看它们是否互为相反数。可是复合函数要连续算两次，直接展开很容易乱。\textbf{不妨只问：第一次计算之后，$2$ 和 $-2$ 变成了什么？} 这正是口诀的来历。

\medskip

\end{explanationbox}
\clearpage
\begin{explanationbox}
\textbf{\textcolor{CExplanation}{二、先从一个初中生也能算的例子开始}}
\begin{center}\begin{tikzpicture}[font=\small,>=stealth,line join=round]
\draw[gray!40,rounded corners=3pt] (0,0) rectangle (14.2,8.3);\node at (7.1,7.700000000000001) {内偶的原因：第一步已经把正负输入合为同一个数};
\node at (7.1,6.9){取 $x=-2,2$，先平方，再加 $3$};\draw[->,gray] (1,5.7)--(12.8,5.7);\node[left] at (1,5.7){$x$};\foreach \X/\v in {2.4/-2,3.5/0,4.6/2}{\draw[gray] (\X,5.7)--(\X,5.6);\node[above=6pt] at (\X,5.7){$\v$};}\fill[blue] (2.4,5.7) circle(2.3pt);\fill[orange!85!black] (4.6,5.7) circle(2.3pt);\draw[->,gray] (1,3.6)--(12.8,3.6);\node[left] at (1,3.6){$u$};\draw[gray] (1.3,3.6)--(1.3,3.5);\node[above=6pt] at (1.3,3.6){$0$};\fill[violet] (3.5,3.6) circle(2.3pt);\node[below left=7pt] at (3.5,3.6){$4$（两路重合）};\draw[->,blue,dashed,shorten >=3pt] (2.4,5.7)--(3.5,3.6);\draw[->,orange!85!black,dashed,shorten >=3pt] (4.6,5.7)--(3.5,3.6);\draw[->,gray] (1,1.8)--(12.8,1.8);\node[left] at (1,1.8){$y$};\draw[gray] (1.3,1.8)--(1.3,1.7);\node[above=6pt] at (1.3,1.8){$0$};\fill[violet] (5.15,1.8) circle(2.3pt);\node[above right=7pt] at (5.15,1.8){$7$（两路仍重合）};\draw[->,violet,dashed,shorten <=3pt,shorten >=3pt] (3.5,3.6)--(5.15,1.8);\node[anchor=west] at (8,4.8){$g(x)=x^2$};\node[anchor=west] at (8,2.8){$f(u)=u+3$};
\node[text width=12.6cm,align=center] at (7.1,.65) {各行单位长度相同，零点位置可不同；所有标数均由实际运算得到。};
\end{tikzpicture}\end{center}


设里面的规则为 $g(x)=x^2$，外面的规则为 $f(u)=u+3$，所以
\[
F(x)=f(g(x))=x^2+3.
\]
拿 $x=2$ 和 $x=-2$ 分别走一遍：
\[
\begin{aligned}
\boxed{2}&\ \xrightarrow{\text{平方}}\ \boxed{4}\ \xrightarrow{\text{加 }3}\ \boxed{7},\\[0.3em]
\boxed{-2}&\ \xrightarrow{\text{平方}}\ \boxed{4}\ \xrightarrow{\text{加 }3}\ \boxed{7}.
\end{aligned}
\]
\textbf{关键不在于“加 3”，而在于“平方”已经让 $2$ 与 $-2$ 得到同一个数 $4$。} 外面即使不是“加 3”，换成开平方、立方或其他运算，只要能算，结果仍然相同。于是整个函数是偶函数。

\medskip

\end{explanationbox}
\clearpage
\begin{explanationbox}
\textbf{\textcolor{CExplanation}{三、内层是奇函数时，为什么要继续看外层？}}
\begin{center}\begin{tikzpicture}[font=\small,>=stealth,line join=round]
\draw[gray!40,rounded corners=3pt] (0,0) rectangle (14.2,8.5);\node at (7.1,7.9) {先看内层：合并同一个数，还是得到一对相反数？};
\node at (2.4,6.9) {$x^2+3$：内偶};\begin{scope}[shift={(2.4,3.8000000000000003)},x=1.1333333333333333cm,y=0.4cm]\draw[->,gray] (-1.5,0)--(1.5,0) node[right]{$x$};\draw[->,gray] (0,-3)--(0,6) node[above]{$y$};\draw[blue,thick,solid,domain=-1.3:1.3,samples=160] plot(\x,{(\x)^2+3});\fill[orange!85!black] (-1,4) circle(2pt);\fill[orange!85!black] (1,4) circle(2pt);\end{scope}\node[align=center,text width=4cm] at (2.4,1.7) {$F(-1)=F(1)=4$};\node at (7.050000000000001,6.9) {$x^3$：内奇外奇};\begin{scope}[shift={(7.050000000000001,3.8000000000000003)},x=1.1333333333333333cm,y=0.4cm]\draw[->,gray] (-1.5,0)--(1.5,0) node[right]{$x$};\draw[->,gray] (0,-3)--(0,6) node[above]{$y$};\draw[violet,thick,solid,domain=-1.3:1.3,samples=160] plot(\x,{(\x)^3});\fill[orange!85!black] (-1,-1) circle(2pt);\fill[orange!85!black] (1,1) circle(2pt);\end{scope}\node[align=center,text width=4cm] at (7.050000000000001,1.7) {$F(-1)=-1,\ F(1)=1$};\node at (11.7,6.9) {$x^2$：内奇外偶};\begin{scope}[shift={(11.7,3.8000000000000003)},x=1.1333333333333333cm,y=0.4cm]\draw[->,gray] (-1.5,0)--(1.5,0) node[right]{$x$};\draw[->,gray] (0,-3)--(0,6) node[above]{$y$};\draw[orange!85!black,thick,solid,domain=-1.3:1.3,samples=160] plot(\x,{(\x)^2});\fill[orange!85!black] (-1,1) circle(2pt);\fill[orange!85!black] (1,1) circle(2pt);\end{scope}\node[align=center,text width=4cm] at (11.7,1.7) {$F(-1)=F(1)=1$};
\node[text width=12.6cm,align=center] at (7.1,.65) {三图同尺度；分别取 $g=x^2,f(u)=u+3$，及 $g=x,f(u)=u^3$ 或 $u^2$。};
\end{tikzpicture}\end{center}


改成里面先“乘 2”，也就是 $g(x)=2x$。这时 $x=1$ 与 $x=-1$ 会先得到 $2$ 与 $-2$，它们不相同，而是互为相反数。

\begin{center}
\renewcommand{\arraystretch}{1.35}
\begin{tabular}{c|c|c|c}
输入 $x$ & 内层 $g(x)=2x$ & 外层规则 & 最后 $F(x)$\\ \hline
$1$ & $2$ & $f(u)=u^3$ & $8$\\
$-1$ & $-2$ & $f(u)=u^3$ & $-8$\\ \hline
$1$ & $2$ & $f(u)=u^2$ & $4$\\
$-1$ & $-2$ & $f(u)=u^2$ & $4$
\end{tabular}
\end{center}
\textbf{看出差别：} 同样是里面“乘 2”，外面用立方，就会保留正负号，整个函数为奇函数；外面用平方，正负号被消去，整个函数为偶函数。\textbf{所以“内奇同外”，不是“内奇一定奇”。}

\medskip

\end{explanationbox}
\clearpage
\begin{explanationbox}
\textbf{\textcolor{CExplanation}{四、用图像认识两种对称}}
\begin{center}\begin{tikzpicture}[font=\small,>=stealth,line join=round]
\draw[gray!40,rounded corners=3pt] (0,0) rectangle (14.2,8.5);\node at (7.1,7.9) {先看内层：合并同一个数，还是得到一对相反数？};
\node at (2.4,6.9) {$x^2+3$：内偶};\begin{scope}[shift={(2.4,3.8000000000000003)},x=1.1333333333333333cm,y=0.4cm]\draw[->,gray] (-1.5,0)--(1.5,0) node[right]{$x$};\draw[->,gray] (0,-3)--(0,6) node[above]{$y$};\draw[blue,thick,solid,domain=-1.3:1.3,samples=160] plot(\x,{(\x)^2+3});\fill[orange!85!black] (-1,4) circle(2pt);\fill[orange!85!black] (1,4) circle(2pt);\end{scope}\node[align=center,text width=4cm] at (2.4,1.7) {$F(-1)=F(1)=4$};\node at (7.050000000000001,6.9) {$x^3$：内奇外奇};\begin{scope}[shift={(7.050000000000001,3.8000000000000003)},x=1.1333333333333333cm,y=0.4cm]\draw[->,gray] (-1.5,0)--(1.5,0) node[right]{$x$};\draw[->,gray] (0,-3)--(0,6) node[above]{$y$};\draw[violet,thick,solid,domain=-1.3:1.3,samples=160] plot(\x,{(\x)^3});\fill[orange!85!black] (-1,-1) circle(2pt);\fill[orange!85!black] (1,1) circle(2pt);\end{scope}\node[align=center,text width=4cm] at (7.050000000000001,1.7) {$F(-1)=-1,\ F(1)=1$};\node at (11.7,6.9) {$x^2$：内奇外偶};\begin{scope}[shift={(11.7,3.8000000000000003)},x=1.1333333333333333cm,y=0.4cm]\draw[->,gray] (-1.5,0)--(1.5,0) node[right]{$x$};\draw[->,gray] (0,-3)--(0,6) node[above]{$y$};\draw[orange!85!black,thick,solid,domain=-1.3:1.3,samples=160] plot(\x,{(\x)^2});\fill[orange!85!black] (-1,1) circle(2pt);\fill[orange!85!black] (1,1) circle(2pt);\end{scope}\node[align=center,text width=4cm] at (11.7,1.7) {$F(-1)=F(1)=1$};
\node[text width=12.6cm,align=center] at (7.1,.65) {三图同尺度；分别取 $g=x^2,f(u)=u+3$，及 $g=x,f(u)=u^3$ 或 $u^2$。};
\end{tikzpicture}\end{center}


偶函数的图像左右一样高：如果 $(a,b)$ 在图像上，$(-a,b)$ 也在图像上，图像关于 $y$ 轴对称。奇函数的图像则成对经过 $(a,b)$ 和 $(-a,-b)$，图像关于原点对称。


这三张图要连起来看：左图的正负输入同高，中图的正负输入对应相反高度；右图虽然长得像一条抛物线，但\textbf{对称轴是 $x=-1$，不是 $y$ 轴}，所以它并不是偶函数。

\medskip

\end{explanationbox}
\clearpage
\begin{explanationbox}
\textbf{\textcolor{CExplanation}{五、一个容易眼前一亮的联系：$\sqrt{x^2}=|x|$}}
\begin{center}\begin{tikzpicture}[font=\small,>=stealth,line join=round]
\draw[gray!40,rounded corners=3pt] (0,0) rectangle (14.2,8.3);\node at (7.1,7.700000000000001) {平方再开方：正负输入都回到同一高度};
\begin{scope}[shift={(4.1,2.307142857142857)},x=1.1538461538461537cm,y=1.3142857142857143cm]\draw[->,gray] (-2.6,0)--(2.6,0) node[right]{$x$};\draw[->,gray] (0,-0.5)--(0,3) node[above]{$y$};\draw[blue,thick,solid,domain=-2.3:2.3,samples=180] plot(\x,{abs(\x)});\fill[orange!85!black] (-1,1) circle(2.2pt);\fill[orange!85!black] (1,1) circle(2.2pt);\end{scope}\node[anchor=west,align=left,text width=5.0cm] at (8.4,4.0) {$F=\sqrt{x^2}=|x|$\\[10pt]$g=x^2\ge0$\\[10pt]都能代入 $f(u)=\sqrt u$\\[10pt]不用要求外层是偶函数};
\node[text width=12.6cm,align=center] at (7.1,.65) {根号函数只接收非负输入；复合之后的定义域却可以是全部实数。};
\end{tikzpicture}\end{center}


取 $g(x)=x^2$，$f(u)=\sqrt{u}$（只能对 $u\geqslant0$ 计算）。那么
\[
F(x)=f(g(x))=\sqrt{x^2}=|x|.
\]
虽然外面的 $\sqrt{u}$ 只在非负数上有定义，不能说它是偶函数，但\textbf{里面先平方，正负数就变成同一个非负数}，所以 $|x|$ 必为偶函数。这说明：\textbf{内偶则偶，并不要求外函数本身具有奇偶性。}

\medskip

\end{explanationbox}
\clearpage
\begin{explanationbox}
\textbf{\textcolor{CExplanation}{六、还有一个出人意料的例子：两层都不对称，整体却对称}}
\begin{center}\begin{tikzpicture}[font=\small,>=stealth,line join=round]
\draw[gray!40,rounded corners=3pt] (0,0) rectangle (14.2,8.3);\node at (7.1,7.700000000000001) {两次偏移抵消，整体也能出现奇偶性};
\begin{scope}[shift={(4.1,3.9499999999999997)},x=1.3043478260869565cm,y=1.0cm]\draw[->,gray] (-2.3,0)--(2.3,0) node[right]{$x$};\draw[->,gray] (0,-2.3)--(0,2.3) node[above]{$y$};\draw[blue,thick,solid,domain=-2:2,samples=180] plot(\x,{\x});\end{scope}\node[anchor=west,align=left,text width=5.0cm] at (8.4,4.0) {$g=x+1$，$f(u)=u-1$\\[10pt]两层本身都非奇非偶\\[10pt]$F=(x+1)-1=x$\\[10pt]整体却为奇函数};
\node[text width=12.6cm,align=center] at (7.1,.65) {口诀给出足够的条件，未穷尽所有复合函数的对称方式。};
\end{tikzpicture}\end{center}


取 $g(x)=x+1$，再取 $f(u)=u-1$。这两个函数都不是奇函数或偶函数，但
\[
F(x)=f(g(x))=(x+1)-1=x
\]
却是奇函数。因为\textbf{外面可以把里面造成的“偏移”再抵消掉}。所以口诀是\textbf{快速判断的可靠办法}，并不是列出了所有可能的情况。

\medskip

\textbf{\textcolor{CExplanation}{七、解题时按什么顺序看？}}
先检查 $F(x)$ 的定义域能否成对取 $x$ 与 $-x$；再计算 $g(-x)$。如果等于 $g(x)$，可直接判偶；如果等于 $-g(x)$，再看外层 $f(-u)$ 与 $f(u)$ 的关系。两种都不符合，就不要硬套口诀，回到 $F(-x)$ 直接比较。
\end{explanationbox}
